\chapter{Methodology}
\label{ch:methodology}

This chapter describes how the decision-support framework works and why it is built the way it
is. It is organised around the sequence in which a planning question passes through the system:
the data that enters it (\Cref{sec:data-provenance,sec:order-heterogeneity}), the time base on
which it is simulated (\Cref{sec:operating-time}), the simulation engine itself
(\Cref{sec:des-engine}), the three sequencing policies that can drive it
(\Cref{sec:policies,sec:rl3}), the capacity search that proposes workforce configurations
(\Cref{sec:capacity-estimate,sec:candidates}), the criteria by which a configuration is
accepted (\Cref{sec:feasibility,sec:economics}), the diagnostics that explain the outcome
(\Cref{sec:stage-pressure,sec:adaptive}), and the experimental design under which everything is
compared (\Cref{sec:experimental-design}). The chapter closes with the verification and
validation performed (\Cref{sec:verification}) and the methodological limitations that follow
from these choices (\Cref{sec:method-limitations}).

Throughout, the description is of the system \emph{as implemented}. Where an implementation
choice departs from a textbook convention, the implemented behaviour is documented and its
consequences discussed rather than silently replaced by the convention.

\section{Data and provenance}
\label{sec:data-provenance}

\subsection{What the framework consumes and what these experiments used}

The framework has two entry points. In \emph{retrospective} mode it consumes an order-level
dataset: one row per order, with an arrival timestamp, an urgency class, an item count and
product attributes. In \emph{forecast} mode it consumes an aggregate demand assumption --- an
expected annual or monthly order volume and an uncertainty level --- and generates plausible
individual orders consistent with it.

It is important to separate the capability from the evidence. The framework \emph{can} consume
a real order-level history; the experiments reported in this thesis \emph{did not}. All data
used in this study are synthetic, but the two workflows do not consume the same synthetic
material, and the distinction matters when reading \Cref{ch:results}.

The \textbf{retrospective} experiments use a synthetic annual dataset of \num{240000}
order-level records covering a full calendar year, constructed specifically for this study.
Every order in that dataset carries the scenario tag \code{seasonal\_base}, and the monthly
order counts reproduce the configured seasonal shares exactly --- January is
\(0.130 \times \num{240000} = \num{31200}\) orders and December is
\(0.170 \times \num{240000} = \num{40800}\) --- which confirms the dataset as generator output
rather than observed data. Each month of that dataset is replayed once, so the retrospective
results rest on a single demand realisation per month.

The \textbf{forecast} experiment does \emph{not} replay that dataset. It starts from an
aggregate December volume assumption and the configured demand profile, and generates fresh
stochastic order streams --- one per replication --- through the same family, complexity and
workload logic as the annual generator (\Cref{sec:future-uncertainty}). Those streams are
therefore statistically consistent with the retrospective dataset without being drawn from it:
the three December replications average \num{41687} orders against the annual dataset's
\num{40800} for the same month, because the sampled monthly volume is itself jittered.

Both workflows are thus synthetic, but only the retrospective one uses the \num{240000}-order
annual dataset. No real, proprietary or client data were used at any point. This is a deliberate delimitation
rather than a convenience: the purpose of the experiments is to characterise how the framework
behaves across a controlled range of demand conditions, and a synthetic generator makes those
conditions explicit and reproducible. The corresponding cost --- that the model has not been
calibrated against a real operation --- is stated among the limitations in
\Cref{sec:method-limitations} and revisited in \Cref{sec:external-validity}.

\subsection{Demand structure}

The generator produces a strongly seasonal year, summarised in \Cref{tab:demand-profile} and
plotted in \Cref{fig:demand}. Demand peaks in the winter campaign months and falls to roughly a
fifth of that level in the summer. Crucially, the urgent share is not constant: it rises with
volume, from around \SI{8}{\percent} in July and August to over \SI{25}{\percent} in December.
The months that are hardest to staff are therefore also the months in which the demanding
service class is most prevalent --- a compounding effect that is central to the results.

\begin{figure}[htbp]
\centering
\includegraphics[width=\textwidth]{f1_demand_profile.pdf}
\caption[Monthly order volume and urgency mix]{Monthly order volume and urgency mix in the
dataset used for the retrospective experiments. Volume and urgent share peak together, so the
months with the greatest capacity requirement are also those with the largest proportion of
tightly constrained orders.}
\label{fig:demand}
\end{figure}

\section{Order heterogeneity and the workload model}
\label{sec:order-heterogeneity}

A central modelling choice is that an order is not a unit of work. Each order carries a
\emph{workload unit} value for each of the three stages, computed when the order is created
from its item count and its product attributes:

\begin{equation}
u_{i,s} \;=\; n_i \times \phi_s(\text{family}_i) \times \gamma_s(\text{complexity}_i)
              \times \omega_s(\text{urgency}_i),
\label{eq:workload-units}
\end{equation}

\noindent where \(u_{i,s}\) is the workload of order \(i\) at stage \(s\), \(n_i\) is its item
count, and \(\phi_s\), \(\gamma_s\) and \(\omega_s\) are stage-specific multipliers for product
family, complexity level and urgency respectively.

The multipliers are what make the model interesting, because they differ across stages. A
\emph{fragile} order carries a multiplier of \num{1.1} at picking but \num{1.8} at packing:
fragile goods are not much harder to retrieve, but they are considerably harder to pack safely.
A \emph{bulky} order is relatively harder at picking (\num{1.3}) than the fragile one, but less
demanding at packing (\num{1.6}) than fragile. Urgency multiplies workload only at dispatch
(\num{1.3} for urgent orders, \num{1.0} otherwise), reflecting the additional handling that
expedited shipments require at the point of release. The full multiplier set is reproduced in
\Cref{app:B}.

The consequence is that the three stages do not see proportional workload. Across the dataset
the mean workload per order is \num{3.39} units at picking, \num{2.46} at packing and
\num{1.17} at dispatch --- a ratio of roughly \(1 : 0.72 : 0.34\) --- and, as
\Cref{fig:stage-workload} shows, that ratio itself shifts through the year as the product and
complexity mix changes. A capacity model that treated an order as one unit of work at every
stage would mis-size at least two of the three stages, and would mis-size them differently in
different months.

\begin{figure}[htbp]
\centering
\includegraphics[width=\textwidth]{f2_stage_workload.pdf}
\caption[Mean workload units per order by stage]{Mean workload units per order at each stage,
by month. The three stages are loaded very differently by the same order stream, and the
relationship between them changes seasonally as the product mix shifts.}
\label{fig:stage-workload}
\end{figure}

\subsection{Service times}

Workload units are converted into service time at execution:

\begin{equation}
t_{i,s} \;=\; \max\!\left(t^{\min}_s,\; \left(b_s + c_s\, u_{i,s}\right) \cdot \varepsilon\right),
\qquad
\varepsilon \sim \operatorname{clip}\!\left(\mathcal{N}(1,\,0.12),\; 0.80,\; 1.25\right),
\label{eq:service-time}
\end{equation}

\noindent where \(b_s\) is a fixed set-up component, \(c_s\) the marginal minutes per workload
unit, and \(\varepsilon\) a clipped multiplicative noise term. The parameters are
\(b_s = \SI{0.5}{\minute}\) at all three stages, with
\(c_{\text{pick}} = 0.90\), \(c_{\text{pack}} = 0.70\) and \(c_{\text{disp}} = 0.65\) minutes
per unit, and \(t^{\min}_s = \SI{0.2}{\minute}\).

Two assumptions are embedded here. The fixed component \(b_s\) represents the per-order
overhead that does not scale with size, so that service time is affine rather than proportional
in workload. The clipped normal noise represents execution variability while ruling out the
extreme draws that an unbounded distribution would occasionally produce; the clipping bounds
(\num{0.80} to \num{1.25}) keep variability within a range plausible for repetitive manual
work. The distribution is an assumption of the model, not an empirical fit --- no real handling
time data were available to fit against, and this is noted as a limitation.

\section{The operating-time abstraction}
\label{sec:operating-time}

\subsection{The problem it solves}

A monthly capacity plan must reconcile two quantities that are easy to state separately and
easy to get wrong together: the hours a worker is physically available to process work, and
the hours the organisation pays for. If demand is spread across a calendar month of roughly
\SI{744}{\hour} while labour cost is computed against \SI{160}{\hour} of paid capacity, then
each simulated worker receives several hundred hours of unpaid processing capacity, and every
resulting capacity recommendation is optimistic.

The framework resolves this by defining a single \emph{operating horizon}:

\begin{equation}
H \;=\; h \times 60 \;=\; 160 \times 60 \;=\; \SI{9600}{\minute},
\label{eq:horizon}
\end{equation}

\noindent where \(h\) is the paid productive hours per worker per month. One worker unit
therefore represents exactly \SI{9600}{} productive worker-minutes, and this same \(h\) appears
in the labour cost of \Cref{eq:cost}. Physical capacity and paid capacity are the same quantity
by construction.

\subsection{How calendar demand is mapped onto it}

Calendar arrival timestamps are compressed onto the operating horizon by a linear rescale
within each order's own calendar month:

\begin{equation}
\tau_i \;=\; \frac{t_i - t^{\text{start}}_{m(i)}}
                  {t^{\text{end}}_{m(i)} - t^{\text{start}}_{m(i)}} \times H,
\label{eq:compression}
\end{equation}

\noindent where \(t_i\) is the calendar arrival time of order \(i\), \(m(i)\) its month, and
\(\tau_i\) its position on the operating clock. Because the mapping is a monotone linear
rescale within a month, the relative ordering of arrivals and the burstiness of the arrival
pattern are preserved exactly; what is not preserved is literal wall-clock time.

This matters for interpretation. The statement ``the warehouse operates \SI{160}{\hour} per
month'' would be wrong. The correct statement is that \SI{160}{\hour} is the monthly productive
capacity of one full-time equivalent, and that demand is evaluated against that capacity with
its temporal shape intact. \Cref{fig:compression} illustrates the mapping.

\begin{figure}[htbp]
\centering
\begin{tikzpicture}[x=0.88cm, y=1cm]
  % calendar axis
  \draw[thick, draw=black!60] (0,2.2) -- (12,2.2);
  \coordinate (calleft) at (0,2.2);
  \node[note, left=2pt of calleft, anchor=east] {calendar\\month};
  \foreach \x in {0,1,...,12} \draw[draw=black!45] (\x,2.1) -- (\x,2.3);
  \node[font=\scriptsize, below=1pt] at (0,2.1) {day 1};
  \node[font=\scriptsize, below=1pt] at (12,2.1) {day 31};

  % a few arrivals, clustered to show burstiness is preserved
  \foreach \x in {0.6, 1.1, 1.35, 4.2, 6.5, 6.8, 7.0, 9.4, 10.6, 10.9, 11.2}
    \fill[thesisblue] (\x,2.2) circle (1.5pt);

  % operating axis
  \draw[thick, draw=thesisblue] (0,0) -- (12,0);
  \coordinate (opleft) at (0,0);
  \node[note, left=2pt of opleft, anchor=east, text=thesisblue]
       {operating\\horizon $H$};
  \node[font=\scriptsize, below=1pt] at (0,-0.05) {$0$};
  \node[font=\scriptsize, below=1pt] at (12,-0.05) {\SI{9600}{\minute}};
  \foreach \x in {0.6, 1.1, 1.35, 4.2, 6.5, 6.8, 7.0, 9.4, 10.6, 10.9, 11.2}
    \fill[thesisblue] (\x,0) circle (1.5pt);

  % mapping arrows
  \foreach \x in {0.6, 1.35, 4.2, 6.8, 9.4, 11.2}
    \draw[-{Stealth[length=1.6mm]}, draw=black!40] (\x,2.1) -- (\x,0.12);

  \coordinate (midright) at (12,1.1);
  \coordinate (below) at (6,-0.1);
  \node[note, right=0.1cm of midright, anchor=west, text width=1.9cm]
       {linear\\rescale};
  \node[note, below=0.55cm of below, text width=12cm]
       {relative ordering and burstiness preserved exactly $\cdot$ wall-clock time not preserved};
\end{tikzpicture}
\caption[Operating-time compression]{Compression of calendar arrivals onto the operating
horizon. Each order's position within its own calendar month is rescaled linearly onto
\([0, H)\), so the shape of demand --- including bursts --- is retained while the clock is
decoupled from calendar time.}
\label{fig:compression}
\end{figure}

\subsection{Consequences for the service-level clock}
\label{sec:sla-clock}

Because the simulation runs on the operating clock, service-level thresholds are measured in
\emph{productive elapsed minutes}. The configured thresholds are \SI{240}{\minute} for urgent
orders and \SI{1440}{\minute} for normal orders. These should be read as service targets on the
operating-time scale, not as literal promises of four hours and twenty-four hours of wall-clock
time to the customer. Mapping them onto real delivery promises would require the calendar
structure of a specific operation --- its shift pattern, its cut-off times and its carrier
schedule --- which is outside the scope of this study.

\subsection{The finite horizon and backlog}
\label{sec:finite-horizon}

The simulation runs for exactly \(H\) minutes and then stops. It does not continue until all
orders are complete. Orders still queued or in service when the horizon ends are treated as
\emph{backlog}: their system time is undefined, they are counted as service-level failures, and
they are additionally reported separately as an unfinished-order count and a backlog share.

This is a deliberate design choice with a clear rationale. If the simulation ran until the last
order finished, an under-staffed configuration would eventually complete all its work and could
report a respectable service level, simply by processing December's orders into what would in
reality be January. The finite horizon prevents insufficient capacity from being concealed by
unlimited post-period processing, and makes the accumulated backlog visible as a distinct
quantity that a planner can reason about.

\section{The discrete-event simulation engine}
\label{sec:des-engine}

\subsection{Why discrete-event simulation}

The choice of discrete-event simulation follows from the structure of the problem rather than
from methodological preference. The quantities that determine whether a workforce configuration
succeeds --- when queues form, how long orders wait, how congestion at one stage propagates to
the next, how much work remains at the end of the month --- are emergent consequences of
individual arrival and service events interacting with finite capacity. They are not available
in closed form for a system with this combination of features: non-stationary bursty arrivals,
heterogeneous and stage-differentiated service requirements, multiple servers per stage, a
sequencing decision at every stage, and two customer classes with different deadlines.

An analytical treatment would require assumptions --- stationarity, exponential service,
class-independent workload --- that remove exactly the features under study. Discrete-event
simulation makes no such assumptions, at the cost of producing estimates rather than closed-form
results \citep{law2015}. The framework therefore uses both: an analytical estimate as a cheap
screening anchor (\Cref{sec:capacity-estimate}), and simulation as the arbiter.

\subsection{Model structure}

The model is implemented in Python using the SimPy process-based simulation library. Orders
arrive according to their compressed operating-time timestamps, join the picking queue, are
served by one of the picking workers, then pass in turn to packing and dispatch. Each stage has
its own queue and its own worker pool; workers are not shared across stages, and queues are
unbounded, so there is no blocking. \Cref{fig:flow} shows the structure.

The state recorded for each order is the set of timestamps at which it entered and left service
at each stage. From these, all reported metrics are derived: system time, service-level
attainment, per-stage waiting time distributions, per-stage utilisation and busy time. Queue
lengths are tracked separately by an event-driven accumulator that records the exact moments at
which a queue gains or loses an item, giving an exact time-weighted average queue length rather
than a sampled approximation.

Utilisation at a stage is computed as busy worker-minutes over available worker-minutes:

\begin{equation}
\rho_s \;=\; \frac{\sum_{i} t_{i,s}}{w_s \cdot H},
\label{eq:utilisation}
\end{equation}

\noindent where \(w_s\) is the number of workers at stage \(s\). Because a worker cannot be
busy for longer than the horizon, a value above \num{1} would indicate an error rather than a
meaningful result, and the implementation raises a warning if one occurs.

\section{The sequencing policies as implemented}
\label{sec:policies}

Three policies are compared. Two are transparent heuristics and one is learned. Because the
precise semantics of these policies determine how the results in \Cref{ch:results} must be
read, each is described exactly as implemented rather than by its conventional definition.

\subsection{FIFO}

Under FIFO each stage holds a single first-in-first-out queue, and workers take the order that
has been waiting longest at that stage. It is important to note that this is
\emph{stage-local}. At picking, queue order is customer arrival order. At packing and dispatch,
it is the order in which upstream work was completed, which need not coincide with arrival
order because service times differ between orders. FIFO in this model therefore means
``first to become available at this stage is first served here'', not ``first ordered by the
customer is first served everywhere''.

FIFO is included as the neutral baseline: it makes no use of the urgency information attached
to an order, and therefore shows what happens when differentiated service commitments are not
reflected in the operating discipline.

\subsection{Urgent-First}
\label{sec:urgent-first}

Under Urgent-First each stage holds a single priority queue, with urgent orders assigned
priority \num{0} and normal orders priority \num{1}. Workers always take an order of the
highest available priority, so every waiting urgent order is served before any normal order.
The policy is non-pre-emptive: an order already in service is never interrupted.

The within-class behaviour requires a precise statement, because it differs from what the name
would conventionally suggest. The priority queue is implemented as a binary heap over
priority-tagged items whose comparison operator considers \emph{only} the priority field. Two
orders of the same class are therefore never ordered relative to one another, and the sequence
in which equal-priority orders leave the queue is determined by the internal structure of the
heap. It is deterministic given the insertion sequence, but it is neither first-in-first-out
nor sorted by order identifier.

This was confirmed by direct observation rather than inferred from the source: inserting ten
equal-priority items in the order \(1,2,\dots,10\) returns them in the order
\(1, 3, 7, 10, 9, 6, 8, 5, 2, 4\). The same check confirms that the class-level priority
\emph{is} strictly respected --- all urgent orders precede all normal orders. The verification
script is described in \Cref{app:E}.

The correct characterisation of this policy is therefore: \emph{strict two-class priority with
unspecified within-class ordering}. This is not a defect to be corrected but a property of the
baseline that must be carried into the interpretation of the results, and
\Cref{sec:disc-policy} returns to it when accounting for the very large p95 waiting times that
Urgent-First produces for the normal class.

\subsection{RL-3}
\label{sec:rl3}

The third policy replaces the fixed rule with a learned one. Under RL-3 each stage holds
\emph{two} first-in-first-out queues, one per class, and a reinforcement learning agent decides
which of the two to serve whenever a worker becomes free.

\paragraph{Agent and scope.} A single deep Q-network \citep{mnih2015} is shared across all
three stages. It is not three separate agents: the same network parameters make the decision at
picking, packing and dispatch, with the identity of the stage supplied as an input feature.
This design isolates the effect of acting at every stage with one consistent policy.

\paragraph{Decision points.} The agent is consulted only when \emph{both} queues at a stage are
non-empty, that is, only when there is a genuine choice to make. If exactly one class is
waiting, the action is forced and the resulting transition is not stored for training, so the
agent is not trained on decisions it did not make.

\paragraph{Actions.} There are two actions: action \num{0} serves the urgent queue and action
\num{1} serves the normal queue. Within the chosen class, service is first-in-first-out ---
which, as noted above, differs from the Urgent-First baseline's within-class behaviour. The two
engines implement different queue mechanics by design, and the comparison is between the
decision rules as implemented.

\Cref{fig:rl-structure} summarises the decision structure.

\begin{figure}[htbp]
\centering
\begin{tikzpicture}[node distance=0.55cm]
  \node[block, minimum width=10.4cm, minimum height=1.0cm] (state)
       {\textbf{state} --- 16 features, all clipped to $[0,1]$\\[1pt]
        \scriptsize six queue lengths $\cdot$ three WIP counts $\cdot$ elapsed time $\cdot$
        head slack (urgent, normal) $\cdot$ stage id $\cdot$ three capacities};

  \node[stage, below=of state, minimum width=3.6cm, minimum height=1.0cm] (dqn)
       {shared DQN};

  \node[block, below=of dqn, minimum width=10.4cm, minimum height=0.95cm] (act)
       {\textbf{action}\quad 0: serve the urgent queue \quad$|$\quad 1: serve the normal queue};

  \node[endnode, below=1.0cm of act, minimum width=10.4cm] (rew)
       {reward assigned at dispatch completion, applied to every decision taken for that order};

  \draw[flow] (state) -- (dqn);
  \draw[flow] (dqn) -- (act);
  \draw[flow, dashed] (act) -- node[note, right=2pt] {deferred} (rew);

  \node[note, above=0.28cm of state, text width=11.5cm]
       {observed at whichever stage is deciding --- consulted only when both queues there
        are non-empty};
  \node[note, right=0.3cm of dqn, text width=4.4cm, anchor=west]
       {the same parameters act at all\\three stages; the stage enters\\as a feature, not as a
        separate agent};
\end{tikzpicture}
\caption[RL-3 decision structure]{The RL-3 decision structure. A single network observes the
present state at whichever stage is deciding and chooses which class to serve; the reward is
deferred until the order leaves the system.}
\label{fig:rl-structure}
\end{figure}

\paragraph{State.} The state is a 16-dimensional vector, all components clipped to
\([0,1]\). It comprises the six queue lengths (urgent and normal at each of the three stages,
normalised by class-specific reference values), the three work-in-progress counts, the fraction
of the horizon elapsed, the normalised slack of the head order in each of the two queues at the
current stage, the stage identifier, and the three normalised worker counts. The complete
specification, including the exact ordering and normalisation of each feature, is given in
\Cref{app:C}.

Two properties of this state deserve emphasis. First, the worker-count features allow the agent
to condition its behaviour on how much capacity is available, so that it need not treat a
tightly constrained configuration identically to a generous one. Second, and more importantly,
\emph{nothing in the state encodes future arrivals}. The agent observes the current queue
state, the current work in progress, the elapsed fraction of the horizon and the urgency of the
orders at the head of each queue. It has no forecast, no lookahead and no privileged
information. Any interpretation of its behaviour as anticipating future demand would be
unsupported.

\paragraph{Reward.} The reward is deferred until an order completes dispatch, at which point
every decision recorded for that order receives the same value:

\begin{equation}
r_i \;=\;
\begin{cases}
w_{c(i)}, & \text{if the order met its service target,}\\[4pt]
-\,p_{c(i)} \cdot \min\!\left(1,\; \dfrac{L_i}{\sigma_{c(i)} \cdot \kappa}\right), & \text{otherwise,}
\end{cases}
\label{eq:reward}
\end{equation}

\noindent where \(c(i)\) is the order's class, \(L_i\) its lateness, \(\sigma_{c(i)}\) its
system-time SLA threshold and \(\kappa\) a cap multiplier. Orders left as backlog at the end of the
horizon receive the full penalty \(-p_{c(i)}\), so that abandoning work is treated as the
service failure it is rather than as a neutral outcome. Coefficient values are given in
\Cref{app:C}.

The structure encodes the intended trade-off: meeting an urgent order's target is worth more
than meeting a normal order's, but the penalty for lateness is the same for both classes, so
the agent has an incentive to protect urgent service without licence to abandon normal service
entirely. Attributing the same terminal reward to every decision along an order's path is an
approximation --- it assumes each decision contributed to the eventual outcome --- and is
recorded as such among the limitations.

\paragraph{Training.} The agent was trained for 200 episodes, each sampling one
(month, workforce configuration) pair uniformly from a fixed training pool, with a learning rate
of \num{1e-3}, discount factor \num{0.99}, batch size \num{256}, a replay buffer of \num{200000}
transitions, a target network updated every \num{2000} steps and \(\varepsilon\)-greedy
exploration annealed from \num{1.0} to \num{0.05}. Training deliberately omits the common random
number mechanism described in \Cref{sec:crn}, so that each training episode encounters
independent service-time realisations. The pool covers three representative months --- June,
October and December, chosen to span low, medium and peak demand --- and, for each, seven
workforce configurations drawn from that month's own candidate neighbourhood, giving 21 training
pairs; two further configurations per month were held out from training entirely.
\Cref{sec:app-training-pool} lists the pool exactly and states what follows from it for
generalisation. Full hyperparameters are in \Cref{app:C}.

\paragraph{Evaluation is greedy.} A distinction that matters for reading every RL-3 figure in
this thesis: exploration is a training device only. Every evaluation run --- the retrospective
grid, the forecast replications and the diagnostics --- invokes the network in greedy mode,
taking the action of maximum estimated value with no random component, and records no training
transitions. The \(\varepsilon\)-greedy schedule above applies during training and nowhere else.
No reported RL-3 metric contains randomly chosen exploratory actions.

\section{Analytical capacity estimate}
\label{sec:capacity-estimate}

Before any simulation is run, an analytical estimate provides a starting point for the search.
For each stage, the expected deterministic workload is divided by the effective capacity of one
worker and rounded up:

\begin{equation}
\hat{w}_s \;=\; \max\!\left(1,\; \left\lceil
\frac{\sum_i \left(b_s + c_s\, u_{i,s}\right)}{H \cdot \upsilon}
\right\rceil\right),
\label{eq:capacity-estimate}
\end{equation}

\noindent where \(\upsilon = 0.85\) is a target utilisation and the noise term of
\Cref{eq:service-time} is fixed at its mean of \num{1}. The target utilisation is a deliberate
safety margin: sizing to \SI{100}{\percent} of capacity would guarantee unbounded queueing in
any system with variability \citep{hopp2011}.

This estimate is an \emph{anchor}, not a recommendation. It knows nothing about queueing
dynamics, sequencing or class-specific service requirements, and \Cref{sec:res-capacity} shows
that it systematically differs from what simulation actually recommends. Its role is to place
the search in the right region of the space so that a bounded set of candidates suffices.

\section{Dynamic candidate generation}
\label{sec:candidates}

Around the analytical centre \((\hat{w}_{\text{pick}}, \hat{w}_{\text{pack}},
\hat{w}_{\text{disp}})\), the framework generates a bounded set of 16 distinct workforce
candidates in a fixed, deterministic order: the centre itself; the client's current workforce
if one was supplied; single-stage perturbations of \(\pm 1\); leaner and safer variants scaling
the whole workforce by \(\mp\SI{15}{\percent}\); balanced two-stage combinations; and finally
single-stage perturbations of \(\pm 2\) if the budget of 16 allows. Every stage is clamped to at
least one worker.

This is a \emph{structured neighbourhood search}, and it is important not to overstate it. It
is not a full grid --- a cubic grid over a plausible range would be prohibitively expensive at
these order volumes --- and it is not global optimisation. There is no guarantee that the best
configuration in the whole space lies within the 16 evaluated. The design trades that guarantee
for tractability, on the reasoning that the analytical anchor places the neighbourhood in a
sensible region and that a planner is choosing among operationally plausible configurations
rather than searching an unbounded space.

Configurations are labelled \rgm{sPKD} for single-digit worker counts and \rgm{sP\_K\_D}
otherwise, as described in the front matter.

\section{Service-level feasibility}
\label{sec:feasibility}

A configuration is \emph{feasible} if and only if it satisfies both class-specific service
floors:

\begin{equation}
\text{feasible} \iff
\left(\text{SLA}_{\text{urgent}} \geq 0.95\right) \;\wedge\;
\left(\text{SLA}_{\text{normal}} \geq 0.80\right).
\label{eq:feasibility}
\end{equation}

The conjunction is the essential point. An aggregate service criterion would be satisfied by
configurations that are operationally unacceptable, because the urgent class is a minority of
volume and its failures are therefore diluted. The results provide a direct illustration: in
October at configuration \rgm{s753}, FIFO achieves \SI{81.0}{\percent} aggregate service
attainment --- a figure that sounds tolerable --- while delivering only \SI{24.5}{\percent}
attainment to the urgent class. Any criterion that reported the aggregate alone would have
accepted that configuration.

Where a configuration is infeasible, the magnitude of failure is summarised by a transparent
violation score, the sum of the two shortfalls, each floored at zero.

\subsection{The recommendation rule as implemented}
\label{sec:selection-rule}

Feasibility is what turns the simulated grid into a recommendation, and the rule that does so is
worth stating exactly, because it is lexicographic rather than a plain cost minimisation.

For a given month the framework holds one simulated row per (configuration, policy) pair. It
then applies, in order:

\begin{enumerate}
  \item \textbf{Eligibility.} In forecast mode, only rows carrying a full three-replication
        aggregate are eligible; single-replication screening rows are retained for diagnostic
        display but can never become the recommendation. In retrospective mode every row is
        eligible.
  \item \textbf{Feasibility first.} If any eligible row satisfies both service floors, the
        recommendation is the cheapest such row by total cost. Infeasible rows are excluded
        entirely at this step, however cheap they are.
  \item \textbf{Fallback.} If no eligible row satisfies both floors, the recommendation is the
        row with the smallest service violation, ties broken by lowest total cost. The output is
        then reported as infeasible.
\end{enumerate}

The fallback branch is stated rather than hidden because it changes how such a recommendation
should be read: it is the least-bad configuration among those examined, not a plan that meets
the commitments. In the retrospective run reported in \Cref{ch:results} the fallback was never
reached --- a feasible row existed in all twelve months --- but the framework's behaviour when
it does not is part of the specification.

One reporting distinction follows from this and is easy to conflate. Alongside the
recommendation, the export layer publishes several \emph{descriptive} per-month statistics for
comparison: the cheapest configuration overall, the smallest configuration meeting both floors,
the smallest meeting the urgent floor alone, and the best result for each individual policy. The
first of these is computed without a feasibility filter and is a summary statistic, not the
selection rule. \Cref{sec:app-rules} tabulates them, and \Cref{tab:recommendations} notes which
column it reports.

\section{Economic model}
\label{sec:economics}

Total monthly cost combines labour and the penalty cost of late orders:

\begin{equation}
C \;=\; \underbrace{\left(w_{\text{pick}} + w_{\text{pack}} + w_{\text{disp}}\right) \cdot k \cdot h}_{\text{labour}}
\;+\; \underbrace{n^{\text{late}}_{\text{urgent}} \cdot \pi_{\text{urgent}}
\;+\; n^{\text{late}}_{\text{normal}} \cdot \pi_{\text{normal}}}_{\text{service failure}},
\label{eq:cost}
\end{equation}

\noindent where \(k\) is the hourly labour cost, \(h\) the paid hours per worker-month, and
\(\pi_c\) the penalty per late order of class \(c\). The same \(h\) appears here and in
\Cref{eq:horizon}, which is the point of the operating-time design.

These values are \emph{scenario assumptions}, supplied at run time, not empirically estimated
parameters. This has an important consequence for how results may be compared. The two
experimental runs reported in this thesis were executed with different economic assumptions:
the retrospective run used \(\pi_{\text{urgent}} = 20\), \(\pi_{\text{normal}} = 5\) and
\(k = 15\); the forecast run used \(\pi_{\text{urgent}} = 15\),
\(\pi_{\text{normal}} = 10\) and \(k = 18\). Each run is therefore reported \emph{as run}, and
\textbf{no absolute monetary comparison is drawn between the two modes} anywhere in this
thesis. Within a run, all comparisons share a common parameter set and are valid.

\section{Stage Pressure: the bottleneck diagnostic}
\label{sec:stage-pressure}

Identifying which stage constrains the operation is a prerequisite for explaining a
recommendation. The framework combines four normalised signals into a single comparable score
per stage:

\begin{equation}
P_s \;=\; 0.40\,\rho_s \;+\; 0.25\,\widetilde{q}^{\,95}_s \;+\; 0.20\,\lambda_s
      \;+\; 0.15\,\widetilde{\ell}_s,
\label{eq:pressure}
\end{equation}

\noindent where \(\rho_s\) is utilisation, \(\widetilde{q}^{\,95}_s\) the p95 waiting time,
\(\lambda_s\) the share of late orders' accumulated waiting time attributable to stage \(s\),
and \(\widetilde{\ell}_s\) the time-weighted average queue length.

The normalisation is not uniform across terms, and describing it accurately matters.
Utilisation enters as an absolute ratio clipped to \([0,1]\), so it is comparable across
scenarios. The p95 waiting time and average queue length are normalised \emph{by the maximum
across the three stages}, so they are relative signals: the most-pressured stage always scores
\num{1.0} on those two components regardless of the absolute values. The late-order waiting
share is already a proportion summing to one across stages. A pressure score is therefore
meaningful for \emph{ranking stages within a scenario}, and should not be compared as an
absolute quantity across scenarios.

\subsection{Why these four signals, and why these weights}

The rationale is operational rather than statistical. Utilisation carries the largest weight
because it is the most direct measure of capacity pressure and the quantity most directly
addressed by adding a worker. The p95 waiting time captures the tail of the waiting
distribution, which is what determines service failures --- a stage can have a modest average
wait and still generate late orders through its tail. The late-order waiting share attributes
observed failures to the stage where the waiting actually accumulated, connecting the diagnostic
to the outcome the planner cares about. Average queue length adds evidence of sustained work
accumulation as opposed to transient spikes. The remaining three weights are ordered to reflect
that decreasing directness of evidence.

It must be stated plainly what this score is not. The weights were \emph{not} obtained by
optimisation, statistical calibration or expert elicitation; they are a transparent heuristic
chosen for interpretability. The score does not establish causality: it identifies where
pressure is concentrated, not what causes any individual order to be late. The appropriate
phrasing --- used throughout \Cref{ch:results} --- is that a stage is ``identified as the
primary pressure location under the diagnostic'', not that it causes the failures observed.

\section{Adaptive capacity search}
\label{sec:adaptive}

The bottleneck diagnostic feeds a further step, but not unconditionally. The search is triggered
only when there is a reason to believe additional capacity might help: either the recommended
configuration is infeasible, or its primary pressure stage runs at \SI{85}{\percent} utilisation
or above \emph{and} late orders are already occurring. Where a recommendation is comfortably
feasible at moderate utilisation, no candidate is constructed and no additional simulation is
spent.

When it is triggered, the framework ranks the stages by pressure and constructs a neighbouring
candidate with one additional worker at the most pressured stage. If the second-ranked stage scores within a configured margin of the
first --- \num{0.05} on the pressure scale --- a candidate is constructed for that stage as
well, so that a near-tie in the diagnostic is resolved by simulation rather than by the score
alone. At most two candidates are tested per iteration, and no stage may be raised beyond its
own analytical estimate plus four workers.

Each candidate is re-simulated under all three policies and compared against the incumbent by an
acceptance rule with three branches, evaluated in order:

\begin{enumerate}
  \item if the incumbent is infeasible and the candidate is feasible, accept --- reaching the
        service floors takes precedence over cost;
  \item if both are feasible and the candidate's total cost is lower, accept;
  \item if both are infeasible but the candidate's service violation is smaller, accept it as
        an improvement towards feasibility.
\end{enumerate}

\noindent Otherwise the candidate is rejected and the incumbent stands. Where the rejection
arises because the candidate is no cheaper, the recorded reason is that the added labour cost
was not offset by the reduction in late-order penalties --- which is the outcome observed
throughout \Cref{sec:res-bottleneck}.

The ordering of these branches matters for interpretation, and it is worth being precise about
it. The procedure is \emph{not} a pure cost minimisation. An additional worker that converts an
infeasible configuration into a feasible one is accepted whatever it costs, because an
infeasible plan is not a plan at all. Cost becomes the deciding criterion only once the service
commitments are already satisfied. After each acceptance the bottleneck ranking is recomputed
from the new incumbent, so successive iterations may address different stages. The procedure
runs for at most four iterations and stops when no tested neighbour improves the objective or
when the per-stage ceilings are reached.

The forecast workflow applies the same acceptance rule with an additional screening step: a
candidate is first evaluated on a single replication, and only if it looks promising is it
validated across all replications before the decision is finalised. This prevents a candidate
assessed on one replication from displacing an incumbent that was validated on three.

The managerial content of this step is worth making explicit, because it is easily
misinterpreted. Identifying a bottleneck is often treated as equivalent to recommending
additional capacity there. It is not. An additional full-time equivalent has a definite monthly
cost, and whether it pays for itself depends on how many service failures it actually prevents
--- a quantity that cannot be read off the bottleneck ranking and must be obtained by
simulating the alternative. \Cref{sec:res-bottleneck} shows that in this study the test was
applied in nine months and the additional worker was rejected in every one of them.

As with the candidate search, this is not optimisation. It is a bounded local improvement step
with an explicit, documented acceptance rule. \Cref{fig:adaptive} shows the loop.

\begin{figure}[htbp]
\centering
\begin{tikzpicture}[node distance=0.5cm]
  \node[block, minimum width=6.8cm] (rec) {recommended configuration from the candidate search};
  \node[block, below=of rec, minimum width=6.8cm] (trig)
       {\textbf{trigger?} infeasible, \emph{or} bottleneck utilisation $\geq \SI{85}{\percent}$
        with late orders};
  \node[block, below=of trig, minimum width=6.8cm] (rank)
       {rank stages by pressure $\rightarrow$ add 1 FTE at the top stage\\[-1pt]
        \scriptsize (and at the second stage if within \num{0.05})};
  \node[block, below=of rank, minimum width=6.8cm, fill=thesisblue!7, draw=thesisblue] (sim)
       {re-simulate the candidate under all three policies};
  \node[block, below=of sim, minimum width=6.8cm] (dec)
       {\textbf{accept?} feasibility gained $\rightarrow$ yes;\ \ cheaper and feasible
        $\rightarrow$ yes;\\[-1pt] smaller violation $\rightarrow$ yes;\ \ otherwise no};
  \node[endnode, below=of dec, minimum width=6.8cm] (fin)
       {final recommendation};

  \foreach \a/\b in {rec/trig, trig/rank, rank/sim, sim/dec}
    \draw[flow] (\a) -- (\b);
  \draw[flow] (dec) -- node[note, right=1pt, pos=0.45] {reject} (fin);
  \draw[flow] (trig.east) -- ++(0.7,0) |- node[note, right=1pt, pos=0.25] {no} (fin.east);
  \draw[flow] (dec.west) -- ++(-0.8,0) |- node[note, left=1pt, pos=0.25] {accept:\\re-rank} (rank.west);
\end{tikzpicture}
\caption[Adaptive capacity search]{The adaptive capacity loop. Additional capacity at the
identified bottleneck is proposed, simulated and then accepted or rejected on explicit
criteria; feasibility takes precedence over cost.}
\label{fig:adaptive}
\end{figure}

\section{Experimental design}
\label{sec:experimental-design}

\subsection{Common random numbers}
\label{sec:crn}

Comparing sequencing policies fairly requires that they face the same work. This is not
automatic: because the policies dequeue orders in different sequences, an engine that sampled
service times as orders entered service would advance its random number stream differently
under each policy, and the same nominal seed would produce different service times.

The framework avoids this by pre-computing a service-time map. For a given scenario seed, the
service time of every (order, stage) pair is drawn once, iterating orders in identifier order
and stages in fixed order, before any policy runs. All three policies then look up the same
map. Order \num{4711} therefore requires exactly the same picking minutes under FIFO, under
Urgent-First and under RL-3; what differs between policies is solely \emph{when} that work is
done, which is precisely the effect under study.

This is the classical variance reduction technique of common random numbers
\citep{law2015}. Its benefit here is that differences observed between policies are attributable
to the sequencing decision rather than to sampling noise, which allows meaningful comparison
without large numbers of replications. \Cref{fig:crn} shows the arrangement.

\begin{figure}[htbp]
\centering
\begin{tikzpicture}[node distance=0.55cm and 0.9cm]
  \node[block, minimum width=8.6cm, fill=thesisblue!8, draw=thesisblue] (map)
       {\textbf{Service-time map}\\[-1pt]
        \footnotesize drawn once per scenario seed, for every (order, stage) pair};

  \node[stage, below=1.25cm of map, minimum width=2.5cm] (uf) {Urgent-First};
  \node[stage, left=of uf,  minimum width=2.5cm] (fifo) {FIFO};
  \node[stage, right=of uf, minimum width=2.5cm] (rl) {RL-3};

  \foreach \p in {fifo,uf,rl} \draw[flow] (map) -- (\p);

  \node[note, below=0.06cm of fifo, text width=2.5cm] {different\\dequeue order};
  \node[note, below=0.06cm of uf,   text width=2.5cm] {different\\dequeue order};
  \node[note, below=0.06cm of rl,   text width=2.5cm] {different\\dequeue order};

  \node[note, above=0.28cm of map, text width=11.5cm, text=black!70]
       {orders iterated in identifier order, stages in fixed order ---
        the map depends only on (orders, seed), never on the policy};
  \node[endnode, below=2.1cm of uf, minimum width=10.2cm] (concl)
       {identical service time per order and stage across all three policies};
  \coordinate (bus) at ([yshift=0.5cm]concl.north);
  \draw[draw=black!70, thick] (fifo.south |- bus) -- (rl.south |- bus);
  \foreach \p in {fifo,rl} \draw[draw=black!70, thick] (\p.south) -- (\p.south |- bus);
  \draw[draw=black!70, thick] (uf.south) -- (uf.south |- bus);
  \draw[flow] (uf.south |- bus) -- (concl.north);
\end{tikzpicture}
\caption[Common random numbers design]{The common random number arrangement. Service times are
drawn once per scenario seed, before any policy runs, and looked up by all three. What differs
between policies is only \emph{when} each unit of work is done.}
\label{fig:crn}
\end{figure}

\subsection{The two workflows}

The framework was exercised in both of its modes, and their designs differ.

The \textbf{retrospective} workflow processed all twelve months of the order-level dataset. For
each month it computed the analytical estimate, generated 16 candidates, and evaluated every
candidate under all three policies --- \(12 \times 16 \times 3 = 576\) simulation runs --- with
one scenario seed per month. Its question is counterfactual: \emph{given the demand that
occurred, which workforce configuration and sequencing policy would the framework have
recommended?} It does not claim that the recommended workforce was the one actually used.

The \textbf{forecast} workflow was exercised on December, the peak month. Starting from an
aggregate expected volume and an uncertainty level, it generated three stochastic replications
of the month's orders and applied a two-stage evaluation. Stage~A screened all 16 candidates
under all three policies on the first replication (\(16 \times 3 \times 1 = 48\) runs); Stage~A
then ranked the configurations and retained four finalists, which Stage~B extended to the two
remaining replications (\(4 \times 3 \times 2 = 24\) further runs). The total is 72 simulation
runs rather than the 144 a full three-replication evaluation would have required, a reduction
of \SI{50}{\percent}. Only the finalists' three-replication aggregates are eligible to become
the recommendation. This produces evidence across scenarios --- a mean cost, the spread across
the three replications and the share of replications in which the service floors were met ---
rather than a single point estimate. Three replications describe that spread; they do not
estimate a distribution, and \Cref{sec:res-future} states the consequences for how the reported
p90 cost and replication success rate should be read.

The asymmetry between the two workflows is a property of the experiments as executed and is
carried into the interpretation: the retrospective results cover a full year at one realisation
per month, while the forecast results cover one month at three replications. They answer
different questions and are not pooled.

\subsection{Forecast scenario generation and uncertainty}
\label{sec:future-uncertainty}

Because RQ4 concerns the \emph{difference} between the two uncertainty structures, the forecast
side deserves to be described mechanically rather than in summary. What follows is the
procedure as implemented.

\paragraph{Inputs.} The forecast entry point takes a planning month, an expected annual order
volume (or a direct monthly override), and a named uncertainty level. The expected monthly
volume is the annual figure multiplied by that month's configured share of the year --- for
December, \num{0.170}. The reported run used the \code{standard} level.

\paragraph{What the uncertainty level is.} An uncertainty level is a pair of configured
coefficients of variation, one for volume and one for the within-month arrival pattern:

\begin{center}
\begin{tabular}{lrr}
\toprule
\textbf{Level} & \textbf{Demand CV} & \textbf{Arrival CV} \\
\midrule
\code{low}      & \num{0.05} & \num{0.15} \\
\code{standard} & \num{0.10} & \num{0.30} \\
\code{high}     & \num{0.20} & \num{0.50} \\
\bottomrule
\end{tabular}
\end{center}

These are \emph{scenario assumptions supplied to the model}, not forecast errors estimated from
any observed forecasting record. Nothing in this study calibrates them against how wrong a real
demand forecast turns out to be. They describe how much variation the planner wishes to test
against, and the reliability figures the workflow reports are conditional on that choice. This
is the sense in which the forecast workflow is a robustness instrument rather than a predictive
one, and it is why \Cref{sec:future-work} identifies propagating genuine forecast error as an
extension rather than a refinement.

\paragraph{How a replication is generated.} Replication \(r\) of month \(m\) is seeded
deterministically as \(42 + 10m + 1000r\), so the three December replications are reproducible
and mutually independent draws. Within a replication:

\begin{enumerate}
  \item \textbf{Volume.} The month's order count is drawn from a lognormal distribution whose
        logarithmic parameters are set so that its mean is the expected monthly volume and its
        coefficient of variation is the level's demand CV. Volume therefore varies from
        replication to replication around the forecast rather than being fixed at it.
  \item \textbf{Arrival pattern.} Daily arrival weights start from the configured calendar
        profile, including the campaign-burst pattern that applies to December, and are then
        multiplied by positive lognormal noise of unit mean whose dispersion is the level's
        arrival CV. This adds day-to-day dispersion on top of the seasonal shape without
        altering the intraday profile.
  \item \textbf{Composition.} Each order's urgency class, product family, complexity level and
        item count are drawn from the same month-specific distributions used by the annual
        generator (\Cref{sec:order-heterogeneity}), and its stage workload units are computed by the
        same rules. The forecast and retrospective streams therefore share a structure while
        being distinct realisations.
\end{enumerate}

\paragraph{What is held common and what varies.} Within one replication, all sixteen
configurations and all three policies are evaluated against exactly the same order stream, and
service times are drawn once per (order, stage) from the common random number map of
\Cref{sec:crn}. Comparisons \emph{within} a replication are therefore matched. Across
replications, everything demand-side changes together: the sampled volume, the arrival pattern
and the individual orders. The variation the reported figures express is thus scenario
variation, not measurement noise on a fixed scenario.

\paragraph{How the four finalists are chosen.} After Stage~A, configurations are ranked using
only replication-1 evidence, by a rule that mirrors the recommendation hierarchy of
\Cref{sec:selection-rule}. A configuration at which \emph{any} policy meets both service floors
ranks ahead of every configuration at which none does, and among those it is ranked by the total
cost of its cheapest feasible policy. Configurations with no feasible policy are ranked after
them, by service violation first and cost second. The four highest-ranked configurations are
carried into Stage~B. The rule operates on configurations, not on (configuration, policy) pairs,
so a finalist carries all three of its policies forward.

\section{Verification and validation}
\label{sec:verification}

Following \citet{sargent2008}, verification and validation are treated as distinct activities.

\textbf{Verification} asks whether the model behaves as designed. Several checks support this.
Utilisation is asserted not to exceed one, which would indicate a worker busy longer than the
horizon. Backlog is confirmed to be a genuine unresolved queue rather than a deadlock: the
horizon timeout guarantees termination regardless of how much work remains. The common random
number mechanism was checked by confirming that service-time lookups depend only on order
identifier, stage and seed. The within-class ordering of the Urgent-First priority queue was
established by direct observation rather than assumption, as described in
\Cref{sec:urgent-first}. Repository-level consistency checks and an audit of the learning agent
are documented in \Cref{app:E}.

An important consequence of the verification stance adopted here is that verification does
\emph{not} license redesign. Where two components were found to implement different queue
mechanics --- the Urgent-First baseline using one priority heap per stage, RL-3 using two
first-in-first-out queues per stage --- the difference is documented as a design fact and
carried into interpretation, not eliminated to make the two engines symmetric.

\Cref{tab:verification} sets out the checks performed and what each establishes.

\begin{table}[htbp]
\centering
\caption[Verification checks]{Verification checks performed, and what each establishes.}
\label{tab:verification}
\small
\begin{tabularx}{\textwidth}{lXX}
\toprule
\textbf{Check} & \textbf{Method} & \textbf{Establishes} \\
\midrule
Utilisation bound & assertion that busy minutes do not exceed worker-minutes available &
  no worker processes work outside the paid horizon \\
Termination & horizon timeout rather than a completion event &
  the run always terminates; leftover work is backlog, not deadlock \\
Backlog accounting & undefined system time propagates to a failed service comparison &
  unfinished orders count as failures and are separately reported \\
Common random numbers & service-time lookup keyed only on order, stage and seed &
  all three policies face identical service times \\
Queue measurement & event-driven accumulator at enqueue/dequeue instants &
  exact time-weighted queue length, not a sampled approximation \\
Within-class ordering & direct observation of dequeue order from the priority structure &
  Urgent-First's within-class semantics, established rather than assumed \\
\bottomrule
\end{tabularx}
\end{table}

The last of these deserves emphasis as a matter of method rather than result. The within-class
ordering of the Urgent-First baseline could have been inferred from reading the source, and an
inference would have been wrong: the natural reading of a policy named ``Urgent-First'' operating
on a priority queue is that equal-priority orders are served first-in-first-out. Running the
structure and observing its output produced the correct answer where reading would not have. The
general lesson --- that behavioural properties of a simulation should be established by
observation wherever they bear on the interpretation of results --- is applied throughout this
study.

\textbf{Validation} asks whether the model is credible for its intended purpose. The validation
performed here is internal and structural rather than external and empirical. The model
reproduces qualitative behaviours that queueing theory predicts: waiting times grow sharply as
utilisation approaches one, congestion concentrates at the stage with the highest load ratio,
and additional capacity at a non-constraining stage produces little improvement. Relative
comparisons behave sensibly: policies that protect a class improve that class's attainment,
larger workforces reduce backlog monotonically, and the cost model responds as specified. What
has \emph{not} been done is validation against a real warehouse: no operational data were
available to compare against, no service-time distributions were fitted empirically, and no
practitioner review of the recommendations was conducted. The model is therefore credible as an
instrument for comparing configurations under stated assumptions, and is not established as
predictive of any specific operation.

\section{Methodological limitations}
\label{sec:method-limitations}

The choices described in this chapter carry limitations that should be stated before the
results are read.

\textbf{Data provenance.} All experiments use synthetic data. The demand structure, service
times and workload multipliers are plausible assumptions, not empirical estimates.

\textbf{Absence of warehouse physics.} Travel time, layout, aisle congestion, batching, waves
and inter-stage blocking are not modelled. The model represents capacity and queueing, and
would under-state congestion effects that arise from physical interference between workers.

\textbf{Operating-time abstraction.} Preserving relative arrival structure while discarding
wall-clock time means the service thresholds are not literal customer promises, and results
cannot be read directly as delivery performance.

\textbf{Monthly FTE, not rosters.} The recommendation is a monthly capacity quantity. Whether
that capacity can be arranged into a workable shift pattern is a separate feasibility question
not addressed here.

\textbf{Heuristic bottleneck weights.} The Stage Pressure weights are chosen for
interpretability and are not calibrated. A different weighting could in principle rank stages
differently, although in the results reported here one stage dominates by a wide margin in
every month.

\textbf{Bounded search.} Sixteen candidates per month is a neighbourhood, not the space. No
optimality claim is made.

\textbf{Single realisation per month in the retrospective workflow.} Retrospective results
carry no replication-based uncertainty estimate. Common random numbers make the
\emph{comparison} between policies fair at a given configuration, but they do not quantify how
much a result would vary under a different realisation of demand.

\textbf{Reward attribution in RL-3.} Assigning a single terminal reward to every decision along
an order's path is a credit-assignment approximation.

\textbf{Reproducibility of the learned policy.} As recorded in \Cref{app:E}, the trained network
weights are not retained in the repository. The persisted evaluation outputs are available and
all reinforcement learning results reported here are taken from them, but exact regeneration
requires the trained checkpoint, which was not available in the environment used for final
thesis preparation.

\textbf{Clipped state features saturate at peak capacity and queue lengths.} The learned
policy's inputs are normalised by fixed constants and clipped to \([0,1]\), so worker counts at
or above 20 per stage, and normal-class queues above 500 orders, all map to the same value.
Both occur throughout the December candidate set (\Cref{sec:app-saturation}). This limits the
resolution with which the agent can distinguish large capacities and severe congestion.
